\documentclass{article}

\usepackage{graphicx}
\usepackage[ngerman]{babel}
\usepackage{amsmath}
\usepackage{amssymb}
\usepackage{amsthm}
\usepackage{url}


\newcommand{\R}{\mathbb{R}}
\newcommand{\C}{\mathbb{C}}
\newcommand{\Q}{\mathbb{Q}}
\newcommand{\Z}{\mathbb{Z}}
\newcommand{\N}{\mathbb{N}}

\newtheorem{thm}{Satz}
\newtheorem{cor}[thm]{Korollar}
\newtheorem{lem}[thm]{Lemma}
\newtheorem{prop}[thm]{Proposition}
\newtheorem{defn}[thm]{Definition}

\title{Approximative Darstellung reeller Zahlen}
\date{5.5.2026}
\author{Max Recker\\ \url{max.recker@uni-bielefeld.de}}

\begin{document}

\maketitle

\section{4.1 b-adische Darstellung reeller Zahlen}

\begin{thm}
Sei \(b\in\N\) mit \(b\geq 2\) und \(x\in\R\setminus\{0\}\). Dann existieren eindeutig bestimmte Zahlen
\[
E\in\Z,\qquad \sigma\in\{-1,1\},\qquad z_i\in\{0,1,\dots,b-1\}
\]
für alle \(i\in\N\cup\{0\}\), so dass
\[
\{\,i\in\N\mid z_i\neq b-1\,\}
\text{ ist unendlich},
\qquad z_0\neq 0,
\]
und
\[
x=\sigma b^E\left(\sum_{i=0}^{\infty}z_i b^{-i}\right).
\]
Diese Darstellung heißt die \emph{normalisierte \(b\)-adische Darstellung} von \(x\).
\end{thm}

\subsection*{Beweis der Aussage}

Sei \(x\in\R\setminus\{0\}\). Das Vorzeichen ist eindeutig durch
\[
\sigma=
\begin{cases}
-1,&x<0,\\
1,&x>0
\end{cases}
\]
festgelegt.

\subsubsection*{Normierung der Mantisse}

Für Ziffern \(z_i\in\{0,\dots,b-1\}\) mit \(z_0\neq 0\) gilt \(z_0\geq 1\). Daher folgt
\[
\sum_{i=0}^{\infty}z_i b^{-i}
=
z_0+\sum_{i=1}^{\infty}z_i b^{-i}
\geq z_0\geq 1.
\]
Andererseits gilt wegen \(z_i\leq b-1\)
\[
\sum_{i=0}^{\infty}z_i b^{-i}
\leq
\sum_{i=0}^{\infty}(b-1)b^{-i}
=
(b-1)\frac{1}{1-\frac1b}
=b.
\]
Da zusätzlich unendlich viele Indizes \(i\in\N\) mit \(z_i\neq b-1\) existieren, kann der Wert \(b\) nicht erreicht werden. Also gilt
\[
1\leq \sum_{i=0}^{\infty}z_i b^{-i}<b.
\]

\subsection*{Bestimmung des Exponenten und Rekursion für die Ziffern}

Aus der Normierungsbedingung
\[
1\leq a_0<b,
\qquad
a_0=b^{-E}|x|
\]
folgt
\[
b^E\leq |x|<b^{E+1}.
\]
Damit ist
\[
E=\left\lfloor\log_b |x|\right\rfloor
\]
eindeutig festgelegt. Nun setzen wir
\[
a_0:=b^{-E}|x|
\]
und definieren rekursiv
\[
z_i:=\lfloor a_i\rfloor,
\qquad
a_{i+1}:=b(a_i-z_i).
\]

Da \(1\leq a_0<b\), gilt \(z_0\neq 0\) und \(z_0\in\{1,\dots,b-1\}\). Außerdem folgt aus \(0\leq a_i-z_i<1\), dass
\[
0\leq a_{i+1}<b.
\]
Also gilt für alle \(i\geq 1\):
\[
z_i=\lfloor a_i\rfloor\in\{0,\dots,b-1\}.
\]

\subsection*{Darstellung von \(a_0\) durch Partialsummen und Restglied}

Wir behaupten:
\[
a_0=
\sum_{i=0}^{n}z_i b^{-i}
+
a_{n+1}b^{-n-1}
\qquad
(n\in\N\cup\{0\}).
\]

\subsubsection*{Beweis durch vollständige Induktion}

Für \(n=0\) gilt wegen \(a_1=b(a_0-z_0)\):
\[
z_0+\frac{a_1}{b}
=
z_0+(a_0-z_0)
=a_0.
\]
Damit ist der Induktionsanfang gezeigt.

Sei nun \(n\in\N\) und die Aussage gelte für \(n-1\), also
\[
a_0=
\sum_{i=0}^{n-1}z_i b^{-i}
+
a_n b^{-n}.
\]
Aus \(a_{n+1}=b(a_n-z_n)\) folgt
\[
a_n=z_n+\frac{a_{n+1}}{b}.
\]
Einsetzen liefert
\[
a_0=
\sum_{i=0}^{n-1}z_i b^{-i}
+
\left(z_n+\frac{a_{n+1}}{b}\right)b^{-n}
=
\sum_{i=0}^{n}z_i b^{-i}
+
a_{n+1}b^{-n-1}.
\]
Damit ist der Induktionsschritt gezeigt.

\subsubsection*{Übergang zur unendlichen Darstellung}

Aus
\[
a_0=
\sum_{i=0}^{n}z_i b^{-i}
+
a_{n+1}b^{-n-1}
\]
folgt durch Grenzübergang \(n\to\infty\), dass
\[
a_0=\sum_{i=0}^{\infty}z_i b^{-i},
\]
denn wegen \(0\leq a_{n+1}<b\) gilt
\[
0\leq a_{n+1}b^{-n-1}<b^{-n}\to 0.
\]

\subsection*{Eindeutigkeit der Darstellung}

Die Bedingung, dass unendlich viele Indizes \(i\) mit \(z_i\neq b-1\) existieren, verhindert Darstellungen wie
\[
1{,}0000\ldots=0{,}9999\ldots.
\]
Wäre die Menge
\[
\{\,i\in\N\mid z_i\neq b-1\,\}
\]
endlich, dann gäbe es ein \(n_0\), so dass \(z_i=b-1\) für alle \(i>n_0\). Dann wäre
\[
\sum_{i=n_0+1}^{\infty}(b-1)b^{-i}=b^{-n_0}.
\]
Man könnte also die Ziffer an Stelle \(n_0\) um \(1\) erhöhen und danach nur noch Nullen schreiben, sofern \(z_{n_0}<b-1\). Genau solche Doppeldarstellungen werden durch die Zusatzbedingung ausgeschlossen.

Angenommen, es gäbe zwei verschiedene normalisierte Darstellungen
\[
\sum_{i=0}^{\infty}y_i b^{-i}
=
\sum_{i=0}^{\infty}z_i b^{-i}.
\]
Sei
\[
n:=\min\{i\in\N_0\mid y_i\neq z_i\}.
\]
Ohne Einschränkung sei \(y_n<z_n\). Dann gilt \(y_n\leq z_n-1\). Für den Rest der linken Darstellung erhält man
\[
\sum_{i=n+1}^{\infty}y_i b^{-i}
\leq
\sum_{i=n+1}^{\infty}(b-1)b^{-i}
=
b^{-n}.
\]
Damit folgt
\[
\sum_{i=0}^{\infty}y_i b^{-i}
\leq
\sum_{i=0}^{n-1}z_i b^{-i}
+
(z_n-1)b^{-n}
+
b^{-n}
=
\sum_{i=0}^{n-1}z_i b^{-i}
+
z_n b^{-n}.
\]
Da die rechte Darstellung wegen der Ziffern nach \(n\) mindestens diesen Wert besitzt, gilt zunächst
\[
\sum_{i=0}^{\infty}y_i b^{-i}
\leq
\sum_{i=0}^{\infty}z_i b^{-i}.
\]
Gleichheit könnte nur eintreten, wenn in der linken Darstellung ab Stelle \(n+1\) nur noch Ziffern \(b-1\) auftreten. Dies ist durch die Normalisierungsbedingung ausgeschlossen. Also gilt sogar
\[
\sum_{i=0}^{\infty}y_i b^{-i}
<
\sum_{i=0}^{\infty}z_i b^{-i},
\]
ein Widerspruch. Damit ist die Darstellung eindeutig.

\section{4.2 Maschinenzahlen}

\begin{defn}
Seien \(b,m\in\N\) mit \(b\geq 2\). Eine \(b\)-adische, \(m\)-stellige, normalisierte Gleitkommazahl hat die Form
\[
x=\sigma b^E\left(\sum_{i=0}^{m-1}z_i b^{-i}\right),
\]
wobei \(\sigma\in\{-1,1\}\), \(E\in\Z\), \(z_0\neq 0\) und \(z_i\in\{0,\dots,b-1\}\) gilt.
\end{defn}

Der Grundgedanke lautet: Ein Rechner kann nur endlich viele Stellen speichern. Deshalb arbeitet man mit endlich vielen Mantissenziffern.

\subsection*{Maschinenzahlbereich}

Für die Darstellung von Zahlen im Rechner müssen Werte für \(b,m,E_{\min},E_{\max}\) gewählt werden. Diese Wahl bestimmt den Maschinenzahlbereich
\[
F(b,m,E_{\min},E_{\max}).
\]
Dieser ist definiert durch
\[
F(b,m,E_{\min},E_{\max})
=
\{0\}\cup
\left\{
\sigma b^E
\left(\sum_{i=0}^{m-1}z_i b^{-i}\right)
\;\middle|\;
\begin{array}{l}
\sigma\in\{-1,1\},\\
E_{\min}\leq E\leq E_{\max},\\
z_0\neq 0,\\
z_i\in\{0,\dots,b-1\}
\end{array}
\right\}.
\]
Die Elemente eines solchen endlichen Maschinenzahlbereichs heißen \emph{Maschinenzahlen}.

\subsection*{Der IEEE-754-Standard}

Im Jahr 1985 wurde der IEEE-754-Standard veröffentlicht. Die normalisierten endlichen \texttt{double}-Zahlen entsprechen im Wesentlichen
\[
F_{\text{double}}:=F(2,53,-1022,1023).
\]

\subsection*{Speicherbedarf von \texttt{double}}

Zur Darstellung einer normalisierten \texttt{double}-Zahl genügen 64 Bit: 1 Bit für das Vorzeichen, 11 Bits für den Exponenten und 52 gespeicherte Bits für die Mantisse.

\subsection*{Das sogenannte Hidden Bit}

Bei normalisierten Binärzahlen gilt stets \(z_0=1\). Dieses erste Mantissenbit muss daher nicht explizit gespeichert werden.

\subsection*{Bitaufbau einer normalisierten \texttt{double}-Zahl}

Die 64 Bits zerfallen in
\[
\boxed{\text{1 Bit}}
\qquad
\boxed{\text{11 Bits}}
\qquad
\boxed{\text{52 Bits}}
\]
für Vorzeichen, Exponent und Mantisse.

Dabei codiert das Vorzeichenbit den Wert \(\frac{1-\sigma}{2}\), die Exponentbits codieren den verschobenen Exponenten \(E+1023\), und die Mantissenbits speichern die Ziffern \(z_1,z_2,\dots,z_{52}\).

Reservierte Exponentenwerte werden für \(0\), subnormale Zahlen, \(+\infty\), \(-\infty\) und NaN verwendet.

\subsection*{Der Darstellungsbereich \(\operatorname{range}(F)\)}

Wir bezeichnen mit
\[
\operatorname{range}(F)
:=
\left[
\min\{f\in F\mid f>0\},
\max F
\right]
\]
das kleinste Intervall, in dem alle positiven Maschinenzahlen liegen.

\section{4.3 Rundungen}

\begin{defn}
Sei
\[
F=F(b,m,E_{\min},E_{\max})
\]
ein Maschinenzahlbereich. Eine Rundung auf \(F\) ist eine Abbildung
\[
\operatorname{rd}:\R\to F,
\]
die jeder Zahl \(x\) eine nächstgelegene Maschinenzahl zuordnet, also
\[
|x-\operatorname{rd}(x)|
=
\min_{a\in F}|x-a|.
\]
\end{defn}

\subsection*{Mehrdeutige Rundungen}

Eine Rundung muss nicht eindeutig sein. Liegt eine Zahl genau in der Mitte zwischen zwei Maschinenzahlen, so haben beide denselben Abstand zu dieser Zahl.

\subsection*{Beispiel}

Im Maschinenzahlbereich \(F(10,2,0,2)\) liegt \(12{,}5\) genau zwischen \(12\) und \(13\). Daher könnte man sowohl auf \(12\) als auch auf \(13\) runden.

\subsection*{IEEE-754-Regel: Runden zur geraden Zahl}

Der IEEE-754-Standard verwendet im Zweifelsfall die Regel, dass so gerundet wird, dass die letzte gespeicherte Stelle gerade ist.

\subsection*{Absolute und relative Fehler}

\begin{defn}
Seien \(x,\tilde{x}\in\R\), wobei \(\tilde{x}\) eine Näherung von \(x\) sei. Dann heißt
\[
|x-\tilde{x}|
\]
der \emph{absolute Fehler}. Falls \(x\neq 0\), heißt
\[
\left|\frac{x-\tilde{x}}{x}\right|
\]
der \emph{relative Fehler}.
\end{defn}

\subsection*{Definition der Maschinengenauigkeit}

\begin{defn}
Sei \(F\) ein Maschinenzahlbereich. Die Maschinengenauigkeit von \(F\) wird definiert durch
\[
\varepsilon(F)
:=
\sup\left\{
\left|\frac{x-\operatorname{rd}(x)}{x}\right|
\;\middle|\;
x\in\R,\ |x|\in\operatorname{range}(F)
\right\}.
\]
Dabei bezeichnet \(\operatorname{rd}\) eine Rundung auf den Maschinenzahlbereich \(F\).
\end{defn}

\subsection*{Satz zur Maschinengenauigkeit}

\begin{thm}
Für jeden Maschinenzahlbereich
\[
F=F(b,m,E_{\min},E_{\max})
\qquad (E_{\min}<E_{\max})
\]
gilt
\[
\varepsilon(F)=\frac{1}{1+2b^{m-1}}.
\]
\end{thm}

\subsection*{Interpretation}

Größeres \(m\) bedeutet kleinere Rundungsfehler und damit genaueres Rechnen. Die Basis \(b\) beeinflusst den Abstand benachbarter Maschinenzahlen.

\subsection*{Beweisidee}

Der größte relative Rundungsfehler entsteht genau zwischen zwei benachbarten Maschinenzahlen bei der kleinsten Größenordnung \(b^{E_{\min}}\). Dort beträgt der Abstand benachbarter Zahlen
\[
b^{E_{\min}}b^{-m+1},
\]
also ist der maximale absolute Fehler
\[
\frac12 b^{E_{\min}}b^{-m+1}.
\]
Für
\[
x=b^{E_{\min}}\left(1+\frac12 b^{-m+1}\right)
\]
liegt \(x\) genau in der Mitte zwischen
\[
b^{E_{\min}}
\quad\text{und}\quad
b^{E_{\min}}\left(1+b^{-m+1}\right).
\]
Damit darf man etwa \(\operatorname{rd}(x)=b^{E_{\min}}\) wählen und erhält
\[
\left|\frac{x-\operatorname{rd}(x)}{x}\right|
=
\frac{\frac12 b^{-m+1}}{1+\frac12 b^{-m+1}}
=
\frac{1}{1+2b^{m-1}}.
\]
Folglich gilt
\[
\varepsilon(F)\geq \frac{1}{1+2b^{m-1}}.
\]

Für die umgekehrte Abschätzung sei \(x\in\operatorname{range}(F)\setminus F\). Zwischen den beiden benachbarten Maschinenzahlen
\[
x'=b^E\sum_{i=0}^{m-1}z_i b^{-i},
\qquad
x''=x'+b^E b^{-m+1}
\]
liegt \(x\). Daher gilt
\[
|x-\operatorname{rd}(x)|
\leq
\frac12 b^E b^{-m+1}.
\]
Da wegen der Normierung \(x\geq b^E\) gilt, folgt
\[
\left|\frac{x-\operatorname{rd}(x)}{x}\right|
\leq
\frac{\frac12 b^E b^{-m+1}}
{b^E+\frac12 b^E b^{-m+1}}
=
\frac{1}{1+2b^{m-1}}.
\]
Also gilt auch
\[
\varepsilon(F)\leq \frac{1}{1+2b^{m-1}}.
\]
Zusammen ergibt sich
\[
\varepsilon(F)=\frac{1}{1+2b^{m-1}}.
\]

\begin{defn}
Sei \(F\) ein Maschinenzahlbereich und \(s\in\N\). Eine Maschinenzahl \(f\in F\) besitzt mindestens \(s\) signifikante Stellen in der \(b\)-adischen Gleitkommadarstellung, falls \(f\neq 0\) gilt und für jede Rundung \(\operatorname{rd}\) sowie jede Zahl \(x\in\R\) mit
\[
\operatorname{rd}(x)=f
\]
die Abschätzung
\[
|x-f|
\leq
\frac12 b^{\lfloor\log_b |f|\rfloor+1-s}
\]
erfüllt ist.
\end{defn}

\section{4.4 Maschinenzahlarithmetik}

Sei \(\circ\) eine der vier Operationen
\[
\{+,-,\cdot,/\}.
\]
Für \(x,y\in F\) gilt im Allgemeinen nicht
\[
x\circ y\in F.
\]
Das exakte mathematische Ergebnis ist also oft keine darstellbare Maschinenzahl.

\subsection*{Beispiel}

\[
10^{300}\cdot 10^{300}=10^{600}
\]
kann außerhalb des Wertebereichs liegen.

\subsection*{Ersatzoperation auf dem Rechner}

Darum verwendet ein Computer nicht die exakte Operation, sondern rechnet mit Rundung zurück in den Maschinenzahlbereich. Wir schreiben dafür:
\[
x\odot y\in F.
\]

\subsection*{Annahme 4.12}

Für eine Rundung \(\operatorname{rd}\) zu \(F\) und alle \(x,y\in F\) gelte
\[
x\odot y=\operatorname{rd}(x\circ y).
\]
Korrektes Runden ist daher wesentlich komplizierter.

\subsection*{Beispiel 4.13}

Sei \(F=F(10,2,-5,5)\) sowie
\[
x=4{,}5\cdot 10^1=45,
\qquad
y=1{,}1\cdot 10^0=1{,}1.
\]
Es werde jeweils auf zwei Mantissenstellen gerundet.

\begin{align*}
x\oplus y
&=\operatorname{rd}(x+y)
=\operatorname{rd}(46{,}1)
=4{,}6\cdot 10^1,\\
x\ominus y
&=\operatorname{rd}(x-y)
=\operatorname{rd}(43{,}9)
=4{,}4\cdot 10^1,\\
x\otimes y
&=\operatorname{rd}(x\cdot y)
=\operatorname{rd}(49{,}5)
\in \{\,4{,}9\cdot 10^1,\;5{,}0\cdot 10^1\,\},\\
x\oslash y
&=\operatorname{rd}\!\left(\frac{x}{y}\right)
=\operatorname{rd}(40{,}9090\ldots)
=4{,}1\cdot 10^1.
\end{align*}

Bei IEEE-Rundung zur geraden Endziffer entsteht im Multiplikationsbeispiel \(50\).

\subsection*{Relativer Fehler einer Rechenoperation}

Gemäß Annahme 4.12 gilt
\[
x\odot y=\operatorname{rd}(x\circ y).
\]
Damit erhält man
\[
\left|
\frac{x\circ y-x\odot y}{x\circ y}
\right|
=
\left|
\frac{x\circ y-\operatorname{rd}(x\circ y)}{x\circ y}
\right|.
\]
Falls
\[
|x\circ y|\in\operatorname{range}(F),
\]
also das exakte Ergebnis im darstellbaren Bereich liegt, folgt aus der Definition der Maschinengenauigkeit:
\[
\left|
\frac{x\circ y-x\odot y}{x\circ y}
\right|
\leq \varepsilon(F).
\]

\end{document}
